\documentclass[11pt]{article}
\usepackage[margin=1in]{geometry}
\usepackage{amsmath,amssymb,amsthm}
\usepackage{graphicx}
\usepackage{hyperref}

\newtheorem{theorem}{Theorem}
\newtheorem{corollary}{Corollary}
\newtheorem{definition}{Definition}

\newcommand{\dist}{\operatorname{dist}}
\newcommand{\spanop}{\operatorname{span}}

\title{A Conditional Lift of Convex Normal Tilings to Schauder-Basis Spaces}
\author{Run \texttt{fa\_banach\_001}, lane 1}
\date{2026-06-15}

\begin{document}
\maketitle

\section*{Status}

This is a conditional reduction packet. It proves a complete implication under
an explicitly stated unproved finite-dimensional hypothesis. It does not solve
the open problem of Deville--Jimenez-Sevilla in full.

\section*{Source and Open Problem}

The source paper is:

\begin{quote}
Robert Deville and Mar Jimenez-Sevilla, \emph{Normal and starlike tilings in
separable Banach spaces}, arXiv:2008.10367.
\end{quote}

In the introduction, after proving universal normal starlike tilings in
separable Banach spaces, the authors state the following open problem.

\begin{center}
\includegraphics[width=\textwidth]{figures/open_problem_crop.png}
\end{center}

Transcription of the relevant part: the question whether every separable
Banach space, finite or infinite dimensional, has a $K$-normal convex tiling
with normality constant $K$ independent of the space and norm remains open. The
authors also say that they do not know whether a positive answer in finite
dimensions implies a positive answer in infinite dimensions.

\section*{Definitions}

A tiling $\{T_i\}$ of a Banach space $X$ is $K$-normal if there are $r>0$ and
points $x_i\in T_i$ such that
\[
   B(x_i,r)\subset T_i\subset B(x_i,Kr)
\]
for all $i$. It is convex if every tile is convex.

For a finite-dimensional Banach space $V$, write $\tau(V)$ for the infimum of
all constants $S$ such that $V$ admits an $S$-normal convex tiling.

\section*{Conditional Dependencies}

The main theorem below has a local hypothesis:
\[
   \sup_k \tau(V_k)<\infty
\]
for the initial finite-dimensional spans of a Schauder basis. The source open
problem's finite-dimensional part would imply this hypothesis with a uniform
constant for every basis.

No proof of the finite-dimensional hypothesis is supplied here.

\section*{Idea of the Proof}

Deville--Jimenez-Sevilla build convex ``cylinders'' $C^k_j$ around the
finite-dimensional spaces $V_k$, then refine every cylinder by a starlike
tiling of $V_k$. Their Proposition 2.6 observes that for any fixed finite
number of layers, one may instead slice those layers by convex tilings of
$V_k$; the normality constant then depends on the largest projection norm and
largest finite-dimensional tiling radius used among those layers.

If the basis projections are uniformly bounded and the finite-dimensional
convex tiling constants of all $V_k$ are uniformly bounded, the same refinement
can be made in every layer at once. Convexity is preserved because each refined
tile is an intersection of a convex cylinder with the inverse image of a convex
finite-dimensional tile.

\begin{theorem}
Let $X$ be an infinite-dimensional Banach space with a Schauder basis
$(e_k)$, partial-sum projections $P_k:X\to V_k=\spanop\{e_1,\ldots,e_k\}$, and
\[
   B=\sup_{k\ge 1}\|P_k\|<\infty.
\]
Assume that $S=\sup_{k\ge 1}\tau(V_k)<\infty$. Then $X$ admits a convex normal
tiling.

More precisely, let $a,b,r,\delta$ be the constants from Lemma 2.1 of
Deville--Jimenez-Sevilla and put
\[
   R=a+2b+\frac{8}{\delta}.
\]
Then $X$ has a convex $K$-normal tiling for every
\[
   K>\frac{B}{r}\bigl(R(1+B)+Sr\bigr).
\]
\end{theorem}

\begin{proof}
We use the notation and construction of Deville--Jimenez-Sevilla. Their Step
3 constructs a tiling $\{C^k_j\}$ of $X$ by convex cylinders, indexed with
$k>0\Rightarrow j>0$, and points $x^k_j$ such that
\[
   B(x^k_j+V_k,r)\subset C^k_j\subset B(x^k_j+V_k,R).
\]
The sets $C^k_j$ are convex because they are intersections of inverse images
of convex quotient tiles under linear quotient maps.

Since $(e_k)$ is a Schauder basis, $X=V_k\oplus \ker P_k$. For each cylinder
we choose the representative $x^k_j$ in $\ker P_k$; this is possible because
every coset modulo $V_k$ has a unique representative in $\ker P_k$.

For each $k\ge 1$, choose an $S$-normal convex tiling
$\{D_i^k\}_i$ of $V_k$, scaled so that for suitable points $d_i^k\in D_i^k$,
\[
   B_{V_k}(d_i^k,r)\subset D_i^k\subset B_{V_k}(d_i^k,Sr).
\]
Scaling is legitimate because normality constants are invariant under
homotheties. For $k=0$ put $D^0_0=\{0\}$.

Define
\[
   \widetilde C^0_{0,j}=C^0_j
\]
and, for $k\ge 1$,
\[
   \widetilde C^k_{i,j}=P_k^{-1}(D_i^k)\cap C^k_j.
\]
Each $\widetilde C^k_{i,j}$ is closed and convex.

The family covers $X$: first choose the unique cylinder $C^k_j$ containing a
given point up to tile boundaries, then choose $i$ with $P_kx\in D_i^k$. The
pairwise interior-disjoint property follows similarly. If two refined tiles
come from different cylinders, the source cylinder tiling already has
disjoint interiors. If they share the same cylinder but have different
finite-dimensional pieces, then any open ball contained in their intersection
would have open image under the surjective map $P_k:X\to V_k$, contradicting
that the finite-dimensional tiling $\{D_i^k\}$ has pairwise interior-disjoint
tiles.

It remains to record uniform normality. For $k=0$, the source estimates give
\[
   B(x^0_j,r)\subset C^0_j\subset B(x^0_j,R).
\]
For $k\ge 1$, we claim
\[
   B\left(x^k_j+d_i^k,\frac{r}{B}\right)
       \subset \widetilde C^k_{i,j}
       \subset B\left(x^k_j+d_i^k,R(1+B)+Sr\right).
\]
For the left inclusion, if
$\|x-(x^k_j+d_i^k)\|\le r/B$, then
\[
   \|P_kx-d_i^k\|\le r
\]
because $P_kx^k_j=0$, hence $P_kx\in D_i^k$. Also
$\dist(x-x^k_j,V_k)\le r/B\le r$, so $x\in B(x^k_j+V_k,r)\subset C^k_j$.
Thus $x\in \widetilde C^k_{i,j}$.

For the right inclusion, take $x\in\widetilde C^k_{i,j}$. Since
$x\in C^k_j$, choose $v\in V_k$ with
\[
   \|x-x^k_j-v\|\le R.
\]
Finite-dimensional subspaces are proximinal, so such a nearest point may be
chosen. Because $x^k_j\in\ker P_k$,
\[
   \|v-P_kx\|=\|P_k(v+x^k_j-x)\|\le B R.
\]
Since $P_kx\in D_i^k\subset B_{V_k}(d_i^k,Sr)$, we get
\[
\begin{aligned}
   \|x-(x^k_j+d_i^k)\|
      &\le \|x-x^k_j-v\|+\|v-P_kx\|+\|P_kx-d_i^k\| \\
      &\le R+BR+Sr=R(1+B)+Sr.
\end{aligned}
\]
Combining the inner and outer estimates gives the stated $K$-normality.
\end{proof}

\begin{corollary}
If there is a universal constant $S$ such that every finite-dimensional Banach
space has an $S$-normal convex tiling, then every Banach space with a monotone
Schauder basis has a convex $K$-normal tiling with $K$ depending only on $S$
and on the constants $a,b,r,\delta$ from the source construction.
\end{corollary}

\begin{proof}
For a monotone basis, $B=\sup_k\|P_k\|=1$. The finite-dimensional hypothesis
gives $\tau(V_k)\le S$ for every $k$, so the theorem applies with a constant
independent of the particular space and norm.
\end{proof}

\section*{Limitations and Novelty Check}

This packet does not settle the finite-dimensional convex-normal tiling
problem. It also does not prove the universal infinite-dimensional statement
for arbitrary separable Banach spaces. For non-monotone bases the constant
obtained above depends on the basis projection constant $B$.

The cheap run indexes were searched for arXiv:2008.10367 and the key phrases
``normal and starlike tilings'', ``normal tilings'', ``starlike tilings'',
``convex normal tiling'', and ``Voronoi''. No existing packet or attempt for
this target was found. No external web novelty search was performed for this
conditional reduction.

\section*{References}

\begin{thebibliography}{9}
\bibitem{DJS}
R. Deville and M. Jimenez-Sevilla,
\emph{Normal and starlike tilings in separable Banach spaces},
arXiv:2008.10367.
\end{thebibliography}

\end{document}
